% Modernised reading — fonds Alexandre Grothendieck, shelfmark 29, batch 3
% (pages 41-60).
%
% This is an INTERPRETATION, not a transcription. It re-reads the pages in
% today's notation and vocabulary, states the hypotheses the manuscript leaves
% implicit, and drops the critical apparatus so that the mathematics reads
% straight through. Everything the brackets and underlines carried moves into a
% footnote. In any dispute of substance the transcription governs: whoever
% cites Grothendieck must cite batch-03.fr.tex.
%
% Scope: the folder was read WHOLE before any of this was written — all eleven
% batches, pages 1-216, in one pass — and one file was then written per batch.
%
% This is the hardest batch of the first half. The run of pages 46-57 is fast
% pencil at 228 ppi and the transcription leaves a great deal of it illegible:
% the STATEMENTS carry, the sentences joining them very often do not. This
% reading states what the statements say and does not invent the joins. Where a
% statement is itself incomplete on the page, it is given incomplete and the
% footnote says so — that is the honest reading of those sheets.
%
% Departures from the page, each footnoted where it occurs:
%   - page 41's middle row is not a condition on coefficients (every family
%     satisfies it); it records the ring, not a completion;
%   - Lemme 2's hypothesis « universellement japonais » is weaker than what the
%     modern statement of the finiteness theorem for f_* assumes;
%   - the run is called a finiteness theorem for pi_1 and never mentions pi_1
%     after its title: the bridge, stated here, is that H^1(X,G) is finite for
%     every finite G exactly when pi_1(X) is topologically of finite type.
%
% Pass: Opus 5 (claude-opus-5), 2026-08-29 — first pass, unchecked against the
% pages by a human.
\documentclass[11pt,a4paper]{article}
% Le préambule commun à toutes les transcriptions du fonds Grothendieck.
%
% Il tient en une idée : **l'incertitude se compose, elle ne se résout pas.**
% Une transcription qui lisse un mot douteux a détruit la seule chose pour
% laquelle on la faisait — un lecteur doit pouvoir savoir, à chaque endroit,
% ce qui était sur le feuillet et ce qui a été ajouté. Les six macros qui
% suivent sont donc l'appareil critique, et non de la décoration.
%
% Les mêmes commandes sont reconnues par `scripts/render.mjs`, qui produit la
% vue de lecture du site. Une macro ajoutée ici doit l'être aussi là-bas,
% faute de quoi le rendu échouera bruyamment — ce qui est voulu : un rendu
% silencieusement faux est pire qu'un rendu absent.

\ProvidesPackage{grothendieck}[2026/08/08 Transcriptions du fonds Grothendieck]

\usepackage{amsmath,amssymb,amsthm}
\usepackage{mathrsfs}
\usepackage{stmaryrd}
% Les diagrammes commutatifs sont partout dans ce fonds : c'est la langue
% même de la géométrie algébrique de Grothendieck, et une transcription qui
% les rendrait en prose ne transcrirait rien.
\usepackage{tikz-cd}
% Français par défaut — c'est la langue des feuillets — mais l'anglais est
% chargé aussi : la traduction et le résumé sont des documents anglais, et la
% typographie française (espaces avant les deux-points) y serait une faute.
% Ces documents basculent par \selectlanguage{english} en tête de corps.
\usepackage[english,french]{babel}
\usepackage{fontspec}
\usepackage{geometry}
\geometry{a4paper,margin=25mm,marginparwidth=28mm}
\usepackage{marginnote}
\usepackage{xcolor}
% ulem, not soul. `soul`'s \st and \ul are built on a character-by-character
% scan that XeLaTeX's font machinery breaks: the document fails with an
% undefined control sequence rather than degrading. ulem does the same two jobs
% and is Unicode-engine safe. `normalem` stops it hijacking \emph, which the
% transcriptions use for Grothendieck's own emphasis.
\usepackage[normalem]{ulem}
\usepackage{hyperref}
\hypersetup{colorlinks=true,linkcolor=black,urlcolor=black,pdfborder={0 0 0}}

% --- Métadonnées du lot -----------------------------------------------------
% Reprises telles quelles par le script de rendu, qui les affiche en tête de
% la vue de lecture. Elles fixent ce que le fichier prétend couvrir : sans
% elles, une transcription flotte sans point d'attache dans un fonds de seize
% mille feuillets.

\newcommand{\folder}[1]{\gdef\grothFolder{#1}}
\newcommand{\batch}[1]{\gdef\grothBatch{#1}}
\newcommand{\leaves}[2]{\gdef\grothFirst{#1}\gdef\grothLast{#2}}
\newcommand{\foldertitle}[1]{\gdef\grothFoldertitle{#1}}
\gdef\grothFolder{?}\gdef\grothBatch{?}\gdef\grothFirst{?}\gdef\grothLast{?}\gdef\grothFoldertitle{}

% --- Le résumé --------------------------------------------------------------
%
% Il ouvre la lecture modernisée, sous le titre « Résumé », et il est composé
% en petit. Ce n'est pas un ornement : c'est le seul passage du document qui
% ne soit pas des mathématiques, et un lecteur doit pouvoir le reconnaître
% avant de l'avoir lu — puis le sauter s'il connaît déjà le sujet, ou s'y
% arrêter s'il ne le connaît pas. Le filet qui le suit marque la reprise du
% corps.

\newenvironment{resume}
  {\small\setlength{\parskip}{0.45em}}
  {\par\medskip\hrule\medskip}

% --- Mots-clés ---------------------------------------------------------------
%
% En anglais, en clôture du résumé de la lecture modernisée : le vocabulaire
% moderne sous lequel chercher ce que les feuillets construisent. Le manifeste
% du site (`npm run manifest`) les extrait pour étiqueter la cote entière —
% cette ligne est la source unique des tags, il n'y a pas de fichier à côté.
\newcommand{\keywords}[1]{\par\smallskip\noindent
  {\footnotesize\textsc{Keywords} --- \textit{#1}}\par}

% --- L'appareil critique ----------------------------------------------------

% Le numéro de feuillet, en marge. C'est l'ancre qui permet de régler un
% désaccord sur une formule en regardant *un* feuillet plutôt qu'une chemise
% de six cents. Le site s'en sert aussi pour tourner les pages du fac-similé
% quand on fait défiler la transcription.
\newcommand{\leaf}[1]{%
  \marginnote{\footnotesize\color{gray!70}\textsf{#1}}%
  \label{leaf:#1}\ignorespaces}

% Illisible : on ne devine pas. Un mot inventé qui se lit comme les autres est
% le pire résultat possible.
\newcommand{\ill}{\textcolor{red!60!black}{[\,\ldots\,]}}

% Lecture douteuse : le mot est donné, et signalé comme incertain.
\newcommand{\uncertain}[1]{\uline{#1}}

% Ajout de l'éditeur — une lettre manquante, un mot sous-entendu.
\newcommand{\add}[1]{\textcolor{blue!50!black}{[#1]}}

% Biffé par Grothendieck lui-même. Ce qu'il a rayé fait partie du document :
% une rature dit souvent où la pensée a changé de direction.
\newcommand{\struck}[1]{\sout{#1}}

% Note du transcripteur, sur l'état matériel du feuillet.
\newcommand{\note}[1]{\par\smallskip{\small\color{gray!60!black}\textit{[#1]}}\par\smallskip}

% Note marginale de Grothendieck, distincte du corps du texte.
\newcommand{\marginal}[1]{\par\smallskip\noindent\hspace{1em}%
  {\small\color{gray!60!black}#1}\par\smallskip}

% --- Titre ------------------------------------------------------------------

\AtBeginDocument{%
  \begin{center}
    {\large\bfseries Fonds Alexandre Grothendieck --- cote n\textsuperscript{o}~\grothFolder}\\[2pt]
    {\itshape\grothFoldertitle}\\[4pt]
    {\small feuillets \grothFirst--\grothLast\ (lot \grothBatch)}\\[2pt]
    {\footnotesize\color{gray!60!black}Transcription établie d'après le fac-similé numérisé
     par l'Université de Montpellier.\\
     Les lectures incertaines sont soulignées, les illisibles notées [\,\ldots\,],
     les ajouts entre crochets.}
  \end{center}
  \vspace{1em}}

\endinput


\folder{29}
\batch{3}
\pages{41}{60}
\dating{1959-1969}
\watermark{Édition de démonstration\\interprétation personnelle de l'œuvre}
\foldertitle{Groupe fondamental [Autour de SGA 4 et SGA 7] — lecture modernisée}

\begin{document}

\section*{Résumé}

\begin{resume}
Le lot précédent s'achevait sur une chemise nouvelle, « Finitude (Pbs et
conjectures) », et sur trois énoncés : le groupe fondamental d'un schéma propre
est engendré par un nombre fini d'éléments ; celui d'un schéma de type fini
l'est encore si l'on ne regarde que la partie première à la caractéristique ;
et de même pour la partie modérée. Ces pages sont le programme qui devait les
démontrer.

Ce que le programme demande vaut d'être dit sans formules. Un groupe
fondamental est un objet infini : il enregistre \emph{tous} les revêtements à
la fois. Demander qu'il soit « topologiquement de type fini », c'est demander
qu'un nombre fini de tours suffise à engendrer tous les autres --- que la
variété, si compliquée soit-elle, n'ait qu'un nombre fini de manières
essentiellement différentes d'être parcourue. Sur les nombres complexes, cela
résulte du fait qu'une variété algébrique se triangule avec un nombre fini de
morceaux ; en caractéristique $p$, il n'y a pas de triangulation, et l'énoncé
est un théorème à trouver.

La méthode de ces pages est de le transformer en un énoncé fini. Un revêtement
principal de groupe $G$ n'est rien d'autre qu'un élément de $H^{1}(X, G)$ ;
et un groupe profini est topologiquement de type fini exactement lorsqu'il n'a
qu'un nombre fini de sous-groupes ouverts de chaque indice. Donc : \emph{le
groupe fondamental de $X$ est de type fini si et seulement si $H^{1}(X,G)$ est
fini pour tout groupe fini $G$}. C'est pourquoi la suite de pages intitulée
« Th.\ de finitude relatif pour $\pi_{1}$ » ne parle plus de $\pi_{1}$ après
son titre : elle démontre trois lemmes sur $H^{1}$, et les assemble en une
proposition principale qui ramène le cas général au cas d'un schéma affine
lisse au-dessus d'un ouvert dense d'une base intègre.

Le premier de ces lemmes est le plus intéressant, et c'est celui qu'on
réutiliserait aujourd'hui : si deux revêtements deviennent isomorphes au-dessus
du point générique \emph{géométrique} de la base, ils le deviennent déjà
au-dessus d'un revêtement fini étale de la base. Autrement dit,
l'isomorphisme, qui semblait n'exister qu'après extension à une clôture
algébrique, est en réalité défini sur une extension finie séparable. La note en
marge --- « ne pas confondre avec les résultats de SGA X, où on prend une fibre
générique \emph{vraie} » --- indique exactement où est la finesse.

Deux feuillets, chacun barré d'un long trait de crayon, font autre chose :
ils construisent la catégorie des modules cohérents d'un schéma formel
\emph{modulo} un sous-schéma de définition, comme une catégorie de fractions,
et trouvent qu'elle est celle du complémentaire. C'est le premier passage du
carton sur les schémas formels, et le lot 5 y reviendra pour de bon.

Les noms modernes à chercher sont \emph{théorème de finitude}, \emph{faisceau
constructible}, \emph{catégorie quotient de Serre}, \emph{théorème d'existence
de Grothendieck}, et pour le premier lemme, \emph{schéma des isomorphismes}.

\keywords{finiteness of the fundamental group, nonabelian H1, constructible sheaf, Serre quotient category, Grothendieck existence theorem}
\end{resume}

\pagerange{41}{42}
\section*{Deux complétions d'un croisement normal, et le groupe fondamental local (pages 41 et 42)}

\subsection*{Le feuillet de calcul}

Soit $A$ un anneau noethérien, $x, y \in A$, tels que $A/x$ et $A/y$ soient
réguliers intègres de dimension 1, et supposons $A$ complet pour la topologie
définie par $xy$. C'est le modèle local d'un croisement normal : deux branches
régulières se coupant en un point.

On dispose alors de trois complétions,
\[
  \hat{A}^{(x)} \;=\; k[y][[x]] , \qquad
  \hat{A}^{(y)} \;=\; k[x][[y]] , \qquad
  \hat{A}^{(x,y)} \;=\; k[[x,y]] ,
\]
et $A$ s'envoie dans les deux premières, qui s'envoient toutes deux dans la
troisième. En termes de la région d'indices $(\alpha, \beta)$ des séries
$\sum a_{\alpha\beta}\, x^{\alpha} y^{\beta}$, que la marge de la page
dessine : $\hat{A}^{(x)}$ n'autorise, pour chaque $\alpha$ borné, qu'un nombre
fini de $\beta$ ; $\hat{A}^{(y)}$ l'inverse ; et $\hat{A}^{(x,y)}$ tout.
Le croquis à l'encre, deux axes $\alpha$ et $\beta$ coupés par une courbe,
figure exactement le découpage de cette région.\footnote{La ligne médiane de la
page décrit $A$ par « seul un nombre fini des $a_{\alpha\beta}$, pour
$\alpha + \beta \leqslant N$ donné, sont non nuls ». Ce n'est pas une
condition : il n'y a qu'un nombre fini de couples $(\alpha,\beta)$ avec
$\alpha+\beta \leqslant N$. La ligne enregistre l'anneau $A$ lui-même, entre
les deux complétions, et non une condition qui le caractériserait. Le feuillet
est par ailleurs traversé de traits de crayon gras et d'un griffonnage qui
recouvrent une partie de l'écriture, et l'hypothèse b) elle-même --- le mot
porté en interligne au-dessus de « complet » --- reste indéterminée.}

C'est la situation locale que la page 106 dessinera --- deux branches et le
petit disque $Z = \operatorname{Spec}\hat{\mathcal{O}}_{X,a}$ autour de leur
croisement --- et celle que les pages 87 à 97 traiteront le long d'une fibre
spéciale.

\subsection*{Invariance birationnelle, et le groupe fondamental local}

La page 42 achève la liste des conjectures ouverte page 40. Son point c) est
l'\emph{invariance birationnelle de la ramification modérée}, marqué « Prouvé »
en marge : c'est ce qui donne un sens au problème, puisqu'il permet ensuite
de procéder par normalisation et par Bertini pour se ramener à une courbe munie
d'un diviseur à croisements normaux, où l'on conclut.

Puis s'ouvre une deuxième question, locale cette fois. Soit $X$ un schéma
noethérien --- complet, ou strictement local et « bon » --- et $U$ un ouvert.
On conjecture :

\begin{enumerate}
  \item[a)] $\pi^{t}_{1}(U)$ topologiquement de type fini ;
  \item[b)] $\pi'_{1}(U)$ de présentation finie.
\end{enumerate}

Par la résolution des singularités et l'invariance birationnelle du cas a), on
doit s'en tirer en se ramenant aux énoncés globaux de la page 40. Et pour a),
par descente on se ramène au cas $X$ normal, puis par section hyperplane ---
du moins si $U = X$ privé de l'origine --- au cas $\dim X = 2$, où l'on dispose
du théorème d'Abhyankar.

En marge, deux notes qui sont des notes à soi-même : « vérifier si c'est
\emph{nécessaire} », et « cela doit marcher \dots ».

\pagerange{43}{45}
\section*{Un schéma formel modulo un sous-schéma de définition (pages 43 et 45)}

Deux feuillets, chacun barré d'un long trait de crayon, construisent deux
catégories quotient. Ils sont hors du fil des conjectures et ils annoncent le
lot 5.

Soit $\mathfrak{X}$ un préschéma formel localement noethérien et
$\mathfrak{X}_{0} = V(\mathcal{J})$ un sous-schéma de définition.

\subsection*{Les modules cohérents}

Soit $\mathrm{Coh}(\mathfrak{X})$ la catégorie des modules cohérents sur
$\mathfrak{X}$, et $\mathfrak{S}$ l'ensemble des flèches $u$ dont le noyau et
le conoyau sont tués par une puissance de $\mathcal{J}$. Cet ensemble admet un
calcul de fractions, et l'on pose
\[
  \mathrm{Coh}(\mathfrak{X}/\mathfrak{X}_{0}) \;=\;
  \mathrm{Coh}(\mathfrak{X})\bigl[\mathfrak{S}^{-1}\bigr] .
\]
En pratique --- c'est la remarque de la page --- on prend la sous-catégorie
épaisse des modules cohérents tués par une puissance de $\mathcal{J}$ et l'on
passe au quotient par elle.\footnote{C'est-à-dire, dans le vocabulaire
d'aujourd'hui, le quotient de Serre de $\mathrm{Coh}(\mathfrak{X})$ par la
sous-catégorie de Serre des modules $\mathcal{J}$-de torsion. Le folio 19 du
même fonds contient la théorie générale de ces quotients, sous le nom de
Gabriel ; ici c'est le cas particulier qui sert.}

\subsection*{L'exemple, et ce qu'il calcule}

Soit $X$ propre sur $Y = \operatorname{Spec} A$, $A$ noethérien séparé et
complet pour un idéal $I$, $Y_{0} = V(I)$, $X_{0} = X \times_{Y} Y_{0}$, et
$\mathfrak{X}$ le complété formel de $X$ le long de $X_{0}$. Le théorème
d'existence donne une équivalence
$\mathrm{Coh}(X) \xrightarrow{\ \sim\ } \mathrm{Coh}(\mathfrak{X})$, compatible
avec les sous-catégories des modules à support dans $X_{0}$ :

\begin{tikzcd}
\mathrm{Coh}(\mathfrak{X}) & \mathrm{Coh}(X) \arrow[l, "\approx"'] \\
\mathrm{Coh}(\mathfrak{X})_{\mathrm{nilp}} \arrow[u] & \mathrm{Coh}(X)_{\mathrm{nilp}} \arrow[u]
\end{tikzcd}

d'où, en passant aux quotients,
\[
  \mathrm{Coh}(\mathfrak{X}/\mathfrak{X}_{0}) \;\approx\;
  \mathrm{Coh}(X - X_{0}) .
\]
C'est le point d'exclamation de la page : ce qui reste du schéma formel une
fois qu'on a divisé par ce qui vit sur la fibre spéciale, c'est le
\emph{complémentaire} de cette fibre.

\subsection*{Les revêtements}

La même construction se fait avec les revêtements. Soit $R(\mathfrak{X})$ la
catégorie des préschémas formels finis sur $\mathfrak{X}$, et disons qu'un
$u : \mathfrak{X}' \to \mathfrak{X}''$ est un \emph{isomorphisme modulo
$\mathfrak{X}_{0}$} si l'homomorphisme d'algèbres correspondant
$\mathcal{A}'' \to \mathcal{A}'$ en est un au sens précédent. On vérifie qu'on
obtient un calcul de fractions à droite, et l'on note
$R(\mathfrak{X}/\mathfrak{X}_{0})$ la catégorie quotient.

Sur le même exemple, le théorème de comparaison donne
$R(\mathfrak{X}) \approx R(X)$, donc
\[
  R(\mathfrak{X}/\mathfrak{X}_{0}) \;\approx\; R(X/X_{0}) \;\approx\;
  R(X - X_{0}) ,
\]
avec, en marge, « à vérifier » : c'est le second isomorphisme qui n'est pas
formel, et la page indique qu'il faudrait passer par le \emph{Main Theorem}
global. Cette question --- \emph{quand un revêtement formel est-il
algébrisable} --- est exactement celle que Murre et Grothendieck traiteront en
1969 sous le nom de théorème de Mumford, aux pages 77 à 85 du même
carton.\footnote{Les deux feuillets sont barrés chacun d'un long trait de
crayon. La biffure porte sur la page entière et non sur telle ou telle ligne ;
ce que le trait signifie --- abandon, ou report ailleurs --- n'est pas dit, et
cette lecture ne le devine pas.}

\pagerange{44}{44}
\section*{Finitude en famille (page 44)}

Un troisième point de la chemise, sous le titre « Théorèmes de finitude et
simultanéité », et marqué « c'est fait » en marge. La question est de savoir
comment les groupes fondamentaux des fibres d'un morphisme de type fini
dépendent du point de la base, et l'on veut des énoncés du type
\emph{constructibilité}. Trois formes, calquées sur a), b), c) de la page 40 :

\begin{enumerate}
  \item[a)] majoration \emph{uniforme} du $\pi_{1}(X_{\bar{s}})$ dans les cas
    de finitude, pour $X/S$ propre ;
  \item[b)] majoration uniforme pour les $\pi'_{1}(X_{\bar{s}})$, ou mieux :
    \emph{la structure de $\pi'_{1}(X_{\bar{s}})$ est fonction constructible de
    $s$} ;
  \item[c)] majoration uniforme pour les $\pi^{t}_{1}(X_{\bar{s}})$.
\end{enumerate}

Et l'estimation : « il semble que les méthodes qui résoudront 1) résoudront
également 3) ». La forme b) --- la constructibilité, et non seulement une
borne --- est la plus forte des trois, et c'est celle que le lemme 1 ci-dessous
prépare.

\pagerange{46}{57}
\section*{« Théorème de finitude relatif pour $\pi_{1}$ » (pages 46 à 57)}

Une suite de six feuillets que Grothendieck pagine lui-même I à VI. Elle est
écrite vite et lourdement raturée ; les énoncés se lisent, les phrases qui les
relient souvent pas. Ce qui suit est ce que les énoncés disent.

\subsection*{Le pont entre $H^{1}$ et $\pi_{1}$}

Il n'est pas écrit sur les pages et il faut l'avoir en tête pour les lire. Pour
$X$ connexe pointé en $x$ et $G$ un groupe fini, l'ensemble $H^{1}(X,G)$ des
classes de revêtements principaux de groupe $G$ s'identifie à
\[
  \operatorname{Hom}_{\mathrm{cont}}\bigl(\pi_{1}(X,x),\, G\bigr) \big/
  \operatorname{int}(G) .
\]
Or un groupe profini est topologiquement de type fini si et seulement s'il n'a,
pour chaque entier $n$, qu'un nombre fini de sous-groupes ouverts d'indice $n$.
Donc : \emph{$\pi_{1}(X,x)$ est topologiquement de type fini si et seulement si
$H^{1}(X,G)$ est fini pour tout groupe fini $G$}. C'est pourquoi la suite,
intitulée « théorème de finitude pour $\pi_{1}$ », ne parle que de $H^{1}$.

\subsection*{Lemme 1 : un isomorphisme générique géométrique descend}

Soit $f : X \to Y$ un morphisme de préschémas localement noethériens, $Y$
intègre de point générique $y$, et soient $X'$, $X''$ deux schémas sur $X$ tels
que $X'_{\bar{y}}$ et $X''_{\bar{y}}$ soient isomorphes au-dessus de
$k(\bar{y})$. Alors il existe un revêtement fini étale $Y_{1}$ de $Y$ tel que
$X' \times_{Y} Y_{1}$ et $X'' \times_{Y} Y_{1}$ soient isomorphes au-dessus de
$X_{Y_{1}}$.

Plus précisément : si $\bar{u} : X'_{\bar{y}} \to X''_{\bar{y}}$ est donné et
si $K_{1}$ est la plus petite sous-extension de $k(\bar{y})$ sur laquelle il
est défini, alors $K_{1}$ est \emph{étale} sur $Y$, et sur le $Y_{1}$
correspondant $\bar{u}$ provient d'un morphisme
$X' \times_{Y} Y_{1} \to X'' \times_{Y} Y_{1}$.

La démonstration introduit $\underline{\mathrm{Hom}}_{X}(X', X'')$ --- ou son
sous-schéma des isomorphismes --- pour se ramener au cas $X' = X$, puis invoque
SGA IX 6.2 pour identifier la fibre en $y$ et conclure par un argument de
profondeur.\footnote{Le corps du raisonnement, aux pages 47 et 49, est presque
entièrement illisible à la définition du fac-similé et la transcription le dit.
Ce qui est restitué ici est ce que les phrases portantes autorisent : le
schéma des morphismes, la réduction à $X' = X$, le renvoi à SGA IX 6.2, et
l'isomorphisme final $X_{1} \approx Z$. Aucune étape n'est reconstruite.}

Deux notes accompagnent le lemme, et elles valent l'énoncé. La première rappelle
la suite exacte d'homotopie
\[
  \pi_{1}(X_{\bar{y}}, \xi) \longrightarrow \pi_{1}(X, \xi) \longrightarrow
  \pi_{1}(Y, \zeta) \longrightarrow 0 ,
\]
que le lemme généralise. La seconde avertit : « ne pas confondre avec les
résultats de SGA X, où on prend une fibre générique \emph{vraie} » --- et où il
faut alors supposer $f$ propre. C'est bien la fibre générique
\emph{géométrique} qui intervient ici, et c'est ce qui rend l'énoncé utile en
famille : il dit que l'isomorphisme, qui semblait n'exister qu'après passage à
une clôture algébrique, est déjà défini sur une extension finie séparable.

\subsection*{Lemme 2 : constructibilité de $f_{*}$}

Soit $f : X \to Y$ de type fini, $Y$ localement noethérien et universellement
japonais. Alors, pour tout faisceau d'ensembles constructible $F$ sur $X$,
$f_{*}(F)$ est constructible.

La marge dit « standard », et le lemme n'est pas démontré.\footnote{L'énoncé
moderne correspondant --- la constructibilité de $R f_{*}$ pour $f$ de type
fini --- se démontre sous l'hypothèse que les schémas soient \emph{excellents}
(le théorème de finitude de Deligne, SGA 4$\frac{1}{2}$). « Universellement
japonais » est la condition de Nagata, plus faible : elle ne porte que sur la
finitude des clôtures intégrales. La lecture retient l'hypothèse de la page et
signale l'écart plutôt que de la corriger, parce que l'usage qui en est fait
ci-dessous --- le lemme 3 --- ne demande que le cas noethérien.}

\subsection*{Lemme 3 : $H^{1}$ le long d'une partition}

Soit $X$ un schéma connexe noethérien, $(X_{i})$ une partition finie de $X$ en
parties localement fermées, et $G$ un groupe fini. Alors l'application
\[
  H^{1}(X, G) \longrightarrow \prod_{i} H^{1}(X_{i}, G)
\]
est à \emph{fibres finies}. En particulier, si $H^{1}(X_{i}, G)$ est fini pour
tout $i$, il en est de même de $H^{1}(X, G)$ --- « kif kif », dit la
page.\footnote{Les deux mots sont bien ceux de la page. Ils reviennent au
folio 152 du même carton (« Kif -- kif », pour dire que deux conditions d'une
proposition disent la même chose) et à la page 190 (« grâce à R1 c'est Kif
Kif »). C'est une locution de travail, non une notation.}

La démonstration passe par la théorie du recollement des faisceaux étales, en
utilisant le lemme 2, et par le fait suivant : si $F$ et $G$ sont des faisceaux
étales constructibles sur $X$ noethérien, alors $\operatorname{Hom}(F, G)$ est
un ensemble fini.

\subsection*{La proposition principale}

Soit $S$ un schéma noethérien universellement japonais et $G$ un groupe fini.
Quatre conditions sont mises en regard, dont la page ne laisse lire que la
forme.\footnote{Les conditions (i) à (iv) des pages 53 et 55 sont écrites en
abrégé et lourdement raturées ; la transcription en laisse une part importante
illisible, et notamment les quantificateurs de (ii) et de (iii) et une partie
de (iv). Ce qui suit est la forme des quatre conditions, non leur énoncé
complet ; le lecteur qui en a besoin doit aller à la transcription et, de là, à
la page.}

\begin{enumerate}
  \item[(i)] Pour tout $X$ de type fini sur $S$, $H^{1}(X, G)$ est fini.
  \item[(ii)] La même chose, restreinte aux $X$ affines lisses de type fini
    au-dessus d'un ouvert dense $S'$ d'un sous-schéma affine intègre non vide
    $S_{0}$ de $S$, à fibres géométriques connexes non vides.
  \item[(iii)] La même chose affaiblie : pour de tels $S'$, $X$, et tout ouvert
    non vide $U$ de $X$, il existe un ouvert non vide $U'$ de $U$ tel que
    $H^{1}(U', G)$ soit fini.
  \item[(iv)] Une condition en deux moitiés : a) une condition fibre à fibre,
    pour $s \in S$ et $\bar{X}$ connexe de type fini sur $s$ ; b) une condition
    sur la base seule, la finitude de $H^{1}(S', H)$ pour $S'$ comme en (ii) et
    $H$ un groupe auxiliaire. Lorsque $H$ est commutatif on peut affaiblir b)
    en y remplaçant $H$ par $G$.
\end{enumerate}

De la démonstration, une seule phrase se lit, et elle suffit à répartir le
travail : \emph{(ii) et (iii) résultent aussitôt du lemme 3 ; c'est pour (iv)
qu'il faut le lemme 1}. Le lemme 3 fait le découpage, le lemme 1 fait le
passage de la fibre géométrique générique à la base.

\subsection*{Les deux corollaires}

\textbf{Corollaire 1.} Supposons $G$ d'ordre premier aux caractéristiques
résiduelles. Supposons que pour tout sous-schéma affine intègre non vide
$S_{0}$ de $S$, tout ouvert dense $S'$ de $S_{0}$ et tout sous-groupe $H$ de
$G \times G$, l'ensemble $H^{1}(S', H)$ soit fini. Alors $H^{1}(X, G)$ est fini
pour tout $X$ de type fini sur $S$.

C'est le passage par $G \times G$ qui mérite l'attention : le lemme 1
\emph{compare deux} revêtements, et les couples qu'il fait intervenir ont pour
groupe structural un sous-groupe de $G \times G$. La condition sur la base est
donc plus forte que celle qu'on croirait avoir à imposer.

\textbf{Corollaire 2.} Deux conditions équivalentes : la finitude de
$H^{1}(X, G)$ pour tout $X$ de type fini et tout $G$ premier aux
caractéristiques résiduelles, et la même chose restreinte comme en
(ii).\footnote{L'énoncé du corollaire 2 est en grande partie illisible sur la
page 55, de même que la note qui le suit ; seule la structure --- deux
conditions dont la seconde est la restriction de la première --- se laisse
lire.}

\subsection*{La récurrence sur la dimension}

La page 57 esquisse comment ce qui précède se démontrerait. Récurrence sur
$\dim X$. Par le lemme 3, on se ramène au cas connexe. Si $\dim X = 1$, $X$ est
lisse, et c'est facile : réduction au cas d'Abhyankar, avec un lemme de
Serre--Lang. Si $\dim X \geqslant 2$, on construit une fibration \emph{lisse}
$X \to Y$ avec $\dim Y < \dim X$, et l'on conclut par l'hypothèse de
récurrence.\footnote{C'est la fin de la suite paginée I à VI, et la page 57 est
la plus abîmée du lot ; la marge porte « finir si le \dots », et la remarque
finale n'est pas lisible. La suite reprend au lot 4, pages 61 à 63, avec les
deux corollaires sur les produits.}

\pagerange{59}{60}
\section*{« $\pi_{1}$ et $\mathrm{Pic}^{\tau}$ birationnels » : la première page (pages 59 et 60)}

Une chemise nouvelle s'ouvre, dont le lot 4 contient la suite. Sa première
page --- Grothendieck la numérote 1 --- porte un théorème et sa démonstration
par récurrence sur la dimension.

\textbf{Théorème.} Soit $X$ un préschéma algébrique sur un corps $k$
algébriquement clos. Alors $\pi'_{1}(X)$, limite projective des quotients
discrets de $\pi_{1}(X)$ d'ordre premier à $p$, est topologiquement de type
fini.

La démonstration procède par réductions successives.

Grâce au théorème de van Kampen on peut supposer $X$ affine, puis affine et
\emph{normal} : sinon, on considère $X' \to X$ un revêtement convenable et le
diagramme $X' \times_{X} X' \rightrightarrows X' \to X$, et $\pi_{1}(X)$ est
engendré par $\pi'_{1}(X')$ et par un nombre fini de générateurs correspondant
aux composantes irréductibles de $X' \times_{X} X'$.

Si $\dim X = 0$, c'est trivial. Si $\dim X = 1$, la théorie de la
spécialisation du groupe fondamental ramène à la caractéristique zéro, où l'on
conclut par voie transcendante.

Si $\dim X \geqslant 2$, on considère un revêtement
$\bar{X}^{n} \to \mathbf{P}^{r}_{k}$ et un sous-espace linéaire $L^{r-n+1}$
défini sur une extension algébriquement close $K$ de $k$, et l'on pose
\[
  C \;=\; \bar{X}^{n} \times_{\mathbf{P}^{r}_{k}} L^{r-n+1} ,
\]
courbe normale sur $K$, « générique ». On montre alors que
\[
  \pi_{1}(C) \longrightarrow \pi_{1}(X)
\]
est \emph{surjectif}, d'où la surjectivité de $\pi'_{1}(C) \to \pi'_{1}(X)$, ce
qui conclut puisque le cas des courbes est acquis.

C'est le théorème de Bertini de Zariski qui fournit cette surjectivité, et le
lot 4 l'énonce et le démontre à la page 66.

\subsection*{La question, et un mot anglais}

La page se termine sur une question, qui est celle vers laquelle tout le
dossier des finitudes tend. Soit $X$ normale complète et $Y \subset X$ fermé.
Le quotient de $\pi_{1}(X - Y)$ par le sous-groupe engendré par les groupes
d'inertie \emph{sauvages} aux composantes de codimension 1 de $Y$ est-il
topologiquement de type fini ?\footnote{Le mot \emph{wild} est en anglais sur
la page, et il y reste : c'est la partie $p$-primaire de l'inertie, celle
qu'aucune extension modérée ne contrôle, et le carton n'a pas de mot français
pour elle. C'est exactement ce que les énoncés a), b), c) de la page 40
cherchaient à écarter en passant à $\pi'_{1}$ ou à $\pi^{t}_{1}$ ; la question
demande si on peut le faire sans jeter le reste.}

\end{document}
